\QHeading{第1题　口算。}
\par\noindent


\begin{oralgrid}
\OralItem{1}{15\times6}{90}
\OralItem{2}{5\times13}{65}
\OralItem{3}{480\div12}{40}
\OralItem{4}{78\div2}{39}
\OralItem{5}{6400\div80}{80}
\OralItem{6}{350\div5}{70}
\OralItem{7}{17\times30}{510}
\OralItem{8}{105\times2}{210}
\OralItem{9}{460\div2}{230}
\OralItem{10}{630\div70}{9}
\OralItem{11}{20\times40}{800}
\OralItem{12}{1000\div8}{125}\end{oralgrid}
\par

\QHeading{口算思路补充（对应上表题号）}

\Answer{上面列出得数，下面按小题给出不同计算思路。计数单位法适用于小数，整数题可比较凑整、拆分或商不变等方法。}
\MethodTitle{第1小题：$15\times 6$}

\MethodBlock{口算思路一：按数位拆分计算}{把一个因数拆成整十、整百数与余下部分之和。}{\begin{aligned}[t]&\phantom{={}}15\times 6\\
&=6\times (10+5)\\
&=6\times 10+6\times 5\\
&=60+30\\
&=90\end{aligned}}
\MethodTitle{第2小题：$5\times 13$}

\MethodBlock{口算思路一：按数位拆分计算}{把一个因数拆成整十、整百数与余下部分之和。}{\begin{aligned}[t]&\phantom{={}}5\times 13\\
&=5\times (10+3)\\
&=5\times 10+5\times 3\\
&=50+15\\
&=65\end{aligned}}
\MethodTitle{第3小题：$480\div 12$}

\MethodBlock{口算思路一：利用商不变的性质}{被除数和除数同时乘同一个非零数，商不变。}{\begin{aligned}[t]&\phantom{={}}480\div 12\\
&=(480\times 5)\div (12\times 5)\\
&=2400\div 60\\
&=40\end{aligned}}

\MethodBlock{口算思路二：拆分除数，连续相除}{}{\begin{aligned}[t]&\phantom{={}}480\div 12\\
&=480\div (2\times 6)\\
&=480\div 2\div 6\\
&=240\div 6\\
&=40\end{aligned}}
\MethodTitle{第4小题：$78\div 2$}

\MethodBlock{口算思路一：拆分被除数}{}{\begin{aligned}[t]&\phantom{={}}78\div 2\\
&=(60+18)\div 2\\
&=60\div 2+18\div 2\\
&=30+9\\
&=39\end{aligned}}

\MethodBlock{口算思路二：利用商不变的性质}{被除数和除数同时乘同一个非零数，商不变。}{\begin{aligned}[t]&\phantom{={}}78\div 2\\
&=(78\times 5)\div (2\times 5)\\
&=390\div 10\\
&=39\end{aligned}}
\MethodTitle{第5小题：$6400\div 80$}

\MethodBlock{口算思路一：同时去掉相同个数的0}{}{\begin{aligned}[t]&\phantom{={}}6400\div 80\\
&=(6400\div 10)\div (80\div 10)\\
&=640\div 8\\
&=80\end{aligned}}

\MethodBlock{口算思路二：拆分除数，连续相除}{}{\begin{aligned}[t]&\phantom{={}}6400\div 80\\
&=6400\div (2\times 40)\\
&=6400\div 2\div 40\\
&=3200\div 40\\
&=80\end{aligned}}
\MethodTitle{第6小题：$350\div 5$}

\MethodBlock{口算思路一：利用商不变的性质}{被除数和除数同时乘同一个非零数，商不变。}{\begin{aligned}[t]&\phantom{={}}350\div 5\\
&=(350\times 2)\div (5\times 2)\\
&=700\div 10\\
&=70\end{aligned}}
\MethodTitle{第7小题：$17\times 30$}

\MethodBlock{口算思路一：按数位拆分计算}{把一个因数拆成整十、整百数与余下部分之和。}{\begin{aligned}[t]&\phantom{={}}17\times 30\\
&=30\times (10+7)\\
&=30\times 10+30\times 7\\
&=300+210\\
&=510\end{aligned}}
\MethodTitle{第8小题：$105\times 2$}

\MethodBlock{口算思路一：按数位拆分计算}{把一个因数拆成整十、整百数与余下部分之和。}{\begin{aligned}[t]&\phantom{={}}105\times 2\\
&=2\times (100+5)\\
&=2\times 100+2\times 5\\
&=200+10\\
&=210\end{aligned}}
\MethodTitle{第9小题：$460\div 2$}

\MethodBlock{口算思路一：利用商不变的性质}{被除数和除数同时乘同一个非零数，商不变。}{\begin{aligned}[t]&\phantom{={}}460\div 2\\
&=(460\times 5)\div (2\times 5)\\
&=2300\div 10\\
&=230\end{aligned}}
\MethodTitle{第10小题：$630\div 70$}

\MethodBlock{口算思路一：同时去掉相同个数的0}{}{\begin{aligned}[t]&\phantom{={}}630\div 70\\
&=(630\div 10)\div (70\div 10)\\
&=63\div 7\\
&=9\end{aligned}}

\MethodBlock{口算思路二：拆分除数，连续相除}{}{\begin{aligned}[t]&\phantom{={}}630\div 70\\
&=630\div (2\times 35)\\
&=630\div 2\div 35\\
&=315\div 35\\
&=9\end{aligned}}
\MethodTitle{第12小题：$1000\div 8$}

\MethodBlock{口算思路一：拆分被除数}{}{\begin{aligned}[t]&\phantom{={}}1000\div 8\\
&=(960+40)\div 8\\
&=960\div 8+40\div 8\\
&=120+5\\
&=125\end{aligned}}

\MethodBlock{口算思路二：利用商不变的性质}{被除数和除数同时乘同一个非零数，商不变。}{\begin{aligned}[t]&\phantom{={}}1000\div 8\\
&=(1000\times 5)\div (8\times 5)\\
&=5000\div 40\\
&=125\end{aligned}}


\par\noindent

\QHeading{第2题　脱式计算。}

\CalcBlock{1}{\begin{aligned}[t]&\phantom{={}}\color{black}32\times[(400-340)\div6]\\&\color{answerred}=32\times[60\div6]\\&\color{answerred}=32\times10\\&\color{answerred}=320\end{aligned}}
\par

\CalcBlock{2}{\begin{aligned}[t]&\phantom{={}}\color{black}182\div[560\div(240-160)]\\&\color{answerred}=182\div[560\div80]\\&\color{answerred}=182\div7\\&\color{answerred}=26\end{aligned}}
\CalcBlock{3}{\begin{aligned}[t]&\phantom{={}}\color{black}[360\div(250-242)]\times15\\&\color{answerred}=[360\div8]\times15\\&\color{answerred}=45\times15\\&\color{answerred}=675\end{aligned}}
\CalcBlock{4}{\begin{aligned}[t]&\phantom{={}}\color{black}560\div(80-45)\times4\\&\color{answerred}=560\div35\times4\\&\color{answerred}=16\times4\\&\color{answerred}=64\end{aligned}}
\CalcBlock[简便解法（补充）]{5}{\begin{aligned}[t]&\phantom{={}}\color{black}86\times [200-(65+76)]\\&\color{answerred}=86\times [200-141]\\&\color{answerred}=86\times 59\\&\color{answerred}=86\times (60-1)\\&\color{answerred}=86\times 60-86\times 1\\&\color{answerred}=5160-86\\&\color{answerred}=5074\end{aligned}}
\MethodBlock{按运算顺序计算（基础练习）}{本单元还要练习运算顺序；与上面的简便解法对照。}{\begin{aligned}[t]&\phantom{={}}86\times[200-(65+76)]\\&=86\times[200-141]\\&=86\times59\\&=5074\end{aligned}}

\MethodBlock{方法二：乘法分配律}{把括号中的两部分分别相乘，再相加或相减。}{\begin{aligned}[t]&\phantom{={}}86\times [200-(65+76)]\\
&=86\times 200-86\times (65+76)\\
&=86\times 200-86\times 141\\
&=17200-12126\\
&=5074\end{aligned}}

\CalcBlock[方法一]{6}{\begin{aligned}[t]&\phantom{={}}\color{black}756-(140\times5-320)\\&\color{answerred}=756-(700-320)\\&\color{answerred}=756-380\\&\color{answerred}=376\end{aligned}}
\MethodBlock{方法二：先处理括号或乘除，再去括号后计算}{加号后的括号去掉不变号，减号后的括号去掉要改变各项的加减号。}{\begin{aligned}[t]&\phantom{={}}756-(140\times 5-320)\\
&=756-(700-320)\\
&=756-700+320\\
&=56+320\\
&=376\end{aligned}}

\par\noindent

\QHeading{第3题　解决问题。}

\Problem{①李老师带800元去买地球仪，买了20个地球仪，还剩40元。一个地球仪多少元？}
\par


\MethodTitle{方法一}
\Answer{\Formula{\begin{aligned}[t]&\phantom{={}}(800-40)\div20\\&=760\div20\\&=38\text{（元）}\end{aligned}}\par 答：一个地球仪38元。}

\MethodTitle{方法二：先平均分带的钱，再扣除每个应分的余钱}
\Answer{800元和剩下的40元都按20个平均分。}
\Answer{\Formula{\begin{aligned}[t]&\phantom{={}}800\div 20-40\div 20\\
&=40-2\\
&=38\text{（元）}\end{aligned}}}
\Answer{答：一个地球仪38元。}
\Problem{②某超市购进了6箱卡通水杯，每箱85个。托运时不小心打碎了20个，剩下的水杯重新装箱，每70个装一箱，可以装多少箱？}

\MethodTitle{方法一}
\Answer{\Formula{\begin{aligned}[t]&\phantom{={}}(85\times6-20)\div70\\&=(510-20)\div70\\&=490\div70\\&=7\text{（箱）}\end{aligned}}\par 答：可以装7箱。}

\MethodTitle{方法二：每箱先装70个，再处理多出的部分}
\Answer{原来6箱各取70个装成6箱，多出的每箱15个合起来，减去损坏的20个，再装箱。}
\Answer{\textbf{6箱多出且未损坏的水杯}\par \Formula{\begin{aligned}[t]&\phantom{={}}(85-70)\times 6-20\\
&=15\times 6-20\\
&=90-20\\
&=70\text{（个）}\end{aligned}}}
\Answer{\textbf{这些水杯还能装}\par \Formula{\begin{aligned}[t]&\phantom{={}}70\div 70\\
&=1\text{（箱）}\end{aligned}}}
\Answer{\textbf{一共可装}\par \Formula{\begin{aligned}[t]&\phantom{={}}6+1\\
&=7\text{（箱）}\end{aligned}}}
\Answer{答：可以装7箱。}
